% GENERATED FILE. Edit canonical lessons, then run npm run generate:books.
% canonical-source-hash: fnv1a64:254a6bb3331ba0eb
% canonical-lessons: MR-C280-mala-vatte, MR-C280-mhanun, MR-R280-opinion-and-result

\chapter{What I Think, and So}
\label{ch:mr-what-i-think-and-so}
\hlchaptermodality{280}

\begin{chapteropening}
\textbf{By the end of this chapter:} \emph{I can give an opinion with malā vāṭtaṁ kī and mājhyā mate, and a result with mhaṇūn and tyāmuḷe.}

Complete MR-R280-opinion-and-result: say what you think and what follows from a reason.
\end{chapteropening}

\section[\texorpdfstring{\mr{मला} \mr{वाटतं} \mr{की} … (malā vāṭtaṁ kī …)}{मला वाटतं की … (malā vāṭtaṁ kī …)}]{\mr{मला} \mr{वाटतं} \mr{की} … (malā vāṭtaṁ kī …) — what I think}

\label{lesson:MR-C280-mala-vatte}



\begin{quote}
Say \textquotedblleft{}I feel fine\textquotedblright{}, with the verb that puts you in the dative. (\emph{malā baraṁ vāṭtaṁ}.) Then \textquotedblleft{}Meera says that Marathi is easy\textquotedblright{}. (\emph{mīrā mhaṇte kī marāṭhī sopaṁ āhe}.)

With \textbf{\mr{की}} you said \emph{malā vāṭtaṁ kī cahā āhe}, \textquotedblleft{}I think there is tea\textquotedblright{}. Put a judgment after \textbf{\mr{की}} instead of a fact, and you have an opinion.
\end{quote}

\subsection*{You'll want to know: \mr{मला} \mr{वाटतं} \mr{की} …}
\begin{quote}
\textbf{\mr{मला} \mr{वाटतं} \mr{की} \mr{उद्या} \mr{पाऊस} \mr{पडेल}.}
\end{quote}

\begin{quote}
\emph{malā vāṭtaṁ kī udyā pāūs paḍel.} — \textquotedblleft{}I think it will rain tomorrow.\textquotedblright{}
\end{quote}

Word for word: \textquotedblleft{}to me it seems that tomorrow rain will fall\textquotedblright{}. The thinker stands in the dative, as with \textbf{\mr{आवडणे}}; \textbf{\mr{की}} opens the opinion, and the opinion is a whole sentence. To ask: \textbf{\mr{तुम्हाला} \mr{काय} \mr{वाटतं}?} (\emph{tumhālā kāy vāṭtaṁ?}), \textquotedblleft{}What do you think?\textquotedblright{} In the past: \textbf{\mr{मला} \mr{वाटलं} \mr{की} …} (\emph{malā vāṭlaṁ kī …}), \textquotedblleft{}I thought that …\textquotedblright{}. \textquotedblleft{}I don't think so\textquotedblright{}: \textbf{\mr{मला} \mr{असं} \mr{वाटत} \mr{नाही}} (\emph{malā asaṁ vāṭat nāhī}), with \textbf{\mr{असं}} (\emph{asaṁ}, \textquotedblleft{}so, like this\textquotedblright{}).

\subsection*{You'll want to know: \mr{माझ्या} \mr{मते}, 'in my opinion'}
\begin{quote}
\textbf{\mr{माझ्या} \mr{मते} \mr{हा} \mr{चित्रपट} \mr{चांगला} \mr{आहे}.}
\end{quote}

\begin{quote}
\emph{mājhyā mate hā citrapaṭ cāṅglā āhe.} — \textquotedblleft{}In my opinion this film is good.\textquotedblright{}
\end{quote}

\textbf{\mr{मते}} is \textbf{\mr{मत}}, \textquotedblleft{}an opinion\textquotedblright{}, the word inside \textbf{\mr{मत} \mr{देणे}}, \textquotedblleft{}to vote\textquotedblright{}, with an ending; \textbf{\mr{माझ्या}} is \textbf{\mr{माझा}} bent to fit it. It sits at the front of the sentence, and it sounds firmer than \textbf{\mr{मला} \mr{वाटतं}}.

\subsection*{Your turn}
Say these aloud:

\begin{itemize}
\raggedright
  \item \textquotedblleft{}I think it will rain tomorrow\textquotedblright{} — \emph{malā vāṭtaṁ kī udyā pāūs paḍel}
  \item \textquotedblleft{}What do you think?\textquotedblright{} — \emph{tumhālā kāy vāṭtaṁ?}
  \item \textquotedblleft{}In my opinion this film is good\textquotedblright{} — \emph{mājhyā mate hā citrapaṭ cāṅglā āhe}
\end{itemize}

\begin{tcolorbox}[breakable,colback=teal!4,colframe=teal!35!black,title={Before you move on}]
\textquotedblleft{}I think that\textquotedblright{}? (\textbf{\mr{मला} \mr{वाटतं} \mr{की}}.) \textquotedblleft{}I thought that\textquotedblright{}? (\textbf{\mr{मला} \mr{वाटलं} \mr{की}}.) \textquotedblleft{}In my opinion\textquotedblright{}? (\textbf{\mr{माझ्या} \mr{मते}}.) \textquotedblleft{}What do you think?\textquotedblright{} (\textbf{\mr{तुम्हाला} \mr{काय} \mr{वाटतं}?})
\end{tcolorbox}
\section[\texorpdfstring{\mr{म्हणून} (mhaṇūn)}{म्हणून (mhaṇūn)}]{\mr{म्हणून} (mhaṇūn) — so, and because of that}

\label{lesson:MR-C280-mhanun}



\begin{quote}
Say \textquotedblleft{}because\textquotedblright{}, the word that answers \textbf{\mr{का}}. (\emph{kāraṇ}.) Then \textquotedblleft{}it was raining\textquotedblright{}. (\emph{pāūs paḍat hotā}.)
\end{quote}

\subsection*{You'll want to know: \mr{म्हणून}, 'so'}
\begin{quote}
\textbf{\mr{पाऊस} \mr{पडत} \mr{होता}, \mr{म्हणून} \mr{मी} \mr{घरी} \mr{राहिलो}.}
\end{quote}

\begin{quote}
\emph{pāūs paḍat hotā, mhaṇūn mī gharī rāhilo.} — \textquotedblleft{}It was raining, so I stayed at home.\textquotedblright{}
\end{quote}

\textbf{\mr{कारण}} gives the reason \textbf{after} the event; \textbf{\mr{म्हणून}} gives the result \textbf{after} the reason. Same two facts, opposite order: \textbf{\mr{मी} \mr{घरी} \mr{राहिलो}, \mr{कारण} \mr{पाऊस} \mr{पडत} \mr{होता}} (\emph{mī gharī rāhilo, kāraṇ pāūs paḍat hotā}).

\textbf{\mr{म्हणून}} is built like \textbf{\mr{करून}}: the \textbf{-\mr{ऊन}} form of \textbf{\mr{म्हणणे}}, \textquotedblleft{}to say\textquotedblright{}. Its past has two shapes: \textbf{\mr{तो} \mr{म्हणाला}} (\emph{to mhaṇālā}), \textquotedblleft{}he said\textquotedblright{}, following the speaker, and \textbf{\mr{त्याने} \mr{म्हटलं}} (\emph{tyāne mhaṭlaṁ}), with \textbf{-\mr{ने}}.

\subsection*{You'll want to know: \mr{त्यामुळे}, 'because of that'}
\begin{quote}
\textbf{\mr{काल} \mr{खूप} \mr{पाऊस} \mr{होता}, \mr{त्यामुळे} \mr{आम्ही} \mr{बाहेर} \mr{गेलो} \mr{नाही}.}
\end{quote}

\begin{quote}
\emph{kāl khūp pāūs hotā, tyāmuḷe āmhī bāher gelo nāhī.} — \textquotedblleft{}Yesterday there was a lot of rain, so we didn't go out.\textquotedblright{}
\end{quote}

\textbf{\mr{त्यामुळे}} (\emph{tyāmuḷe}) is \textbf{\mr{त्या}}, \textquotedblleft{}that\textquotedblright{}, and \textbf{-\mr{मुळे}}, \textquotedblleft{}because of\textquotedblright{}. It leans on the cause: \textquotedblleft{}because of that\textquotedblright{}. Where \textbf{\mr{म्हणून}} often marks a choice someone made, \textbf{\mr{त्यामुळे}} marks what came of something.

\subsection*{Your turn}
Say these aloud:

\begin{itemize}
\raggedright
  \item \textquotedblleft{}It was raining, so I stayed at home\textquotedblright{} — \emph{pāūs paḍat hotā, mhaṇūn mī gharī rāhilo}
  \item \textquotedblleft{}I stayed at home, because it was raining\textquotedblright{} — \emph{mī gharī rāhilo, kāraṇ pāūs paḍat hotā}
  \item \textquotedblleft{}Yesterday there was a lot of rain, so we didn't go out\textquotedblright{} — \emph{kāl khūp pāūs hotā, tyāmuḷe āmhī bāher gelo nāhī}
\end{itemize}

\begin{tcolorbox}[breakable,colback=teal!4,colframe=teal!35!black,title={Before you move on}]
Which word comes after the reason? (\textbf{\mr{म्हणून}}.) Which comes before it? (\textbf{\mr{कारण}}.) \textquotedblleft{}Because of that\textquotedblright{}? (\textbf{\mr{त्यामुळे}}.)
\end{tcolorbox}
\section[(review)]{Review — opinion and result}

\label{lesson:MR-R280-opinion-and-result}



\begin{quote}
Say \textquotedblleft{}after school I had tea; after that I went home and slept\textquotedblright{}. (\emph{śāḷenantar mī cahā ghetlā. tyānantar mī gharī jāūn jhoplo.}) Then name the season Arun likes, with \textbf{\mr{ऋतू}}.
\end{quote}

\subsection*{Your turn}
Say these aloud:

\begin{itemize}
\raggedright
  \item I think Meera will come tomorrow.
  \item In my opinion this play is good.
  \item It was raining, so I didn't go to school.
  \item Yesterday there was a lot of rain, so the river was full.
  \item which season you like, and why, with \textbf{\mr{ऋतू}} and \textbf{\mr{कारण}}
\end{itemize}

\begin{tcolorbox}[breakable,colback=teal!4,colframe=teal!35!black,title={Before you move on}]
\textquotedblleft{}I think that\textquotedblright{}? (\textbf{\mr{मला} \mr{वाटतं} \mr{की}}.) \textquotedblleft{}So\textquotedblright{}? (\textbf{\mr{म्हणून}}.) \textquotedblleft{}Because of that\textquotedblright{}? (\textbf{\mr{त्यामुळे}}.) Which letter starts \textbf{\mr{ऋतू}}? (\textbf{\mr{ऋ}}.)
\end{tcolorbox}
